% ---------------------------------------------------------------------------
% MLSys 2027 — main-track submission (skeleton, A-stage of A→B staging)
%
% Status: skeleton + section outlines + framing prose. Body to be filled
% in iteratively. Empirical numbers in §5 (Evaluation) are marked
% [TODO:bench] and depend on the B-stage bench harness landing
% (paper/mlsys-2027/README.md → "A → B staging").
%
% Build: `make` in this directory (pdflatex + bibtex + pdflatex × 2).
% Venue style file (mlsys2027.sty) will be swapped in when the CFP
% drops; until then we compile with the generic article class.
% ---------------------------------------------------------------------------

\documentclass[11pt]{article}

% --- Geometry / typography ---
\usepackage[letterpaper,margin=1in]{geometry}
\usepackage{microtype}
\usepackage[utf8]{inputenc}
\usepackage[T1]{fontenc}
\usepackage{times}

% --- Math ---
\usepackage{amsmath}
\usepackage{amssymb}
\usepackage{amsthm}
\usepackage{mathtools}

% --- Tables / figures ---
\usepackage{graphicx}
\usepackage{booktabs}
\usepackage{multirow}
\usepackage{subcaption}

% --- Code listings ---
\usepackage{listings}
\usepackage{xcolor}
\lstset{
  basicstyle=\ttfamily\small,
  commentstyle=\color{gray!70!black}\itshape,
  keywordstyle=\color{blue!50!black}\bfseries,
  stringstyle=\color{red!60!black},
  showstringspaces=false,
  breaklines=true,
  numbers=left,
  numberstyle=\tiny\color{gray},
  frame=single,
  framerule=0pt,
  backgroundcolor=\color{gray!8},
  xleftmargin=2em,
  xrightmargin=1em,
  belowskip=-0.5em
}
\lstdefinelanguage{cppx}[]{C++}{
  morekeywords={requires,concept,consteval,constinit,co_await,co_return,co_yield},
  moredelim=[s][\color{purple!70!black}]{"}{"_tex},
}

% --- References / hyperlinks ---
\usepackage[hidelinks]{hyperref}
\usepackage{cleveref}
\usepackage{url}

% --- TODO marker (no external package dep) ---
\newcommand{\TODO}[1]{\textbf{\textcolor{red!70!black}{[TODO: #1]}}}
\newcommand{\BENCH}[1]{\textbf{\textcolor{orange!70!black}{[BENCH: #1]}}}

% --- Domain macros (will migrate to paper/common/symbols.tex) ---
\newcommand{\NamedTensor}[1]{\mathsf{#1}}
\newcommand{\ax}[1]{\textsf{#1}}                              % axis name
\newcommand{\udltex}[1]{\texttt{"#1"\_tex}}                   % _tex UDL literal
\newcommand{\code}[1]{\texttt{#1}}

% ---------------------------------------------------------------------------

\title{Named Axes as a Systems Discipline:\\
       Hybrid Compile-Time and Runtime Coordination\\
       for Tensor Programs}

\author{%
  \TODO{author block — fill at submission time}\\
  \texttt{authors@example.com}
}

\date{}

\begin{document}
\maketitle

% ===========================================================================
\begin{abstract}
% ~150-200 words. The "elevator" answer to: what is the contribution,
% what changes if a reader accepts it, what evidence we offer.
Tensor programs routinely fail on \emph{axis-bug} classes that current
type systems cannot catch: silent broadcasts across mis-named dimensions,
transposed contractions, and shape-compatible but semantically incorrect
reductions. Existing named-axis libraries address this either through
runtime string-based notation (\code{einops}, \code{einx}) or
language-level named tensors with fully runtime checks (\code{Haliax},
\code{Penzai}, \code{xarray}). We argue these approaches under-use the
host language's static type system, and present a hybrid design that
encodes axis information across the compile-time / runtime boundary.
The contributions are: (i) a \udltex{i j {->} i j} user-defined-literal
that lifts Einstein-style formulas into C++ template parameters at
compile time, (ii) a hybrid \emph{NTTP $\times$ DynamicShape} axis
system that statically rejects axis mismatches when shapes are known
and falls back to runtime guards otherwise, and (iii) a
\emph{Hexagonal-lite} kernel-backend port that demonstrates substrate
substitutability across a reference implementation, an Eigen-based
substrate, and a WebGPU substrate (Dawn) behind a single port boundary.
\TODO{1-sentence headline empirical result once the B-stage bench
harness lands — e.g., ``We catch 87\% of the axis-bug benchmark suite at
compile time with zero runtime cost; \BENCH{X}-fold speedup on the
bundle-adjustment case study.''}
\end{abstract}

% ===========================================================================
\section{Introduction}
\label{sec:intro}

% ~1 page. Lead with the failure-mode story (axis bug), position
% against the landscape, state contributions, lay out roadmap.

Consider a researcher writing a multi-head attention block. They
compose two linear projections, a softmax, and a matmul. The shapes
type-check at every step --- every intermediate has the right number of
dimensions --- yet the network silently produces incorrect outputs at
training time, because two axes were transposed five lines up. The
type system saw a $(B,H,T,D)$ tensor where the writer meant
$(B,T,H,D)$; the names lived only in a comment.

This failure mode is endemic. \TODO{cite a concrete bug-hunting study
or anecdote --- ideally one of: PyTorch named-tensors postmortem,
Haliax design doc on motivation, our own M1 landscape report
\cite{landscape-comparison-2026}.} Existing libraries respond along
two axes:

\begin{itemize}
  \item \textbf{String-based runtime notation} (\code{einops}~\cite{einops},
        \code{einx}~\cite{einx}): expressive (Einstein-style formulas
        are the source of truth) but defers all checking to runtime
        and forfeits the host language's type system.
  \item \textbf{Language-level named tensors with runtime checks}
        (\code{Haliax}~\cite{haliax}, \code{Penzai}~\cite{penzai},
        \code{xarray}~\cite{xarray}, deprecated
        \code{torch.named\_tensor}~\cite{pytorch-named-tensors}):
        axis names are first-class values, but verification still
        happens at runtime.
\end{itemize}

\paragraph{The thesis.}
A named-axis library should be able to refute the attention-block bug
\emph{at compile time, when the shapes are known}, while still
permitting the runtime-shape case that batch processing demands. The
test of a design is not whether axis names are present but whether they
participate in the type system.

\paragraph{Contributions.}
This paper presents the \code{tensor} library, organised around three
design moves:

\begin{enumerate}
  \item \textbf{The \udltex{} bridge} (\Cref{sec:udl}): a user-defined
        literal that parses Einstein-style formulas at compile time
        and emits the corresponding template parameter pack.
        \udltex{i j {->} i j} \emph{is} the type signature.
  \item \textbf{Hybrid axis system} (\Cref{sec:hybrid}): non-type
        template parameters carry compile-time axis sizes;
        \code{DynamicShape} carries the runtime-determined cases; the
        two compose without bifurcating the kernel API.
  \item \textbf{Hexagonal-lite kernel-backend port}
        (\Cref{sec:hex-lite}): a single port boundary admits three
        substrates (reference, Eigen, WebGPU-via-Dawn~\cite{dawn,adr0014})
        without leaking substrate concerns into the named-axis layer.
\end{enumerate}

\paragraph{Roadmap.}
\Cref{sec:bg} surveys the axis-bug taxonomy and the prior-art landscape.
\Cref{sec:udl,sec:hybrid,sec:hex-lite} present the three design moves.
\Cref{sec:eval} reports the static-catch benchmark and a
bundle-adjustment case study across three backends. \Cref{sec:related}
positions against contemporaneous work. \Cref{sec:discussion} discusses
limitations.

% ===========================================================================
\section{Background \& Motivation}
\label{sec:bg}

% ~1 page. Three subsections: axis-bug taxonomy, prior-art landscape
% (compressed from M1 study), the open gap our design fills.

\subsection{The axis-bug taxonomy}
\label{sec:bg-taxonomy}

% [TODO: concrete enumeration with one-line example each]
% Drawn from internal landscape report; aim for 4-6 classes:
%   T1. Transposed contraction (axes match by position, not meaning)
%   T2. Silent broadcast (size-1 dim gets stretched against wrong axis)
%   T3. Reduction over wrong axis (sum-over-batch vs sum-over-feature)
%   T4. Shape-correct but semantically swapped (B,H vs H,B)
%   T5. Concatenation along the wrong axis
%   T6. Stacking when concatenation was meant (or vice versa)

\TODO{taxonomy table with 4--6 axis-bug classes, each with a 1-line
\code{torch}/\code{numpy} reproduction. Cross-reference against the
M1 landscape report.}

\subsection{Landscape}
\label{sec:bg-landscape}

% Compressed from M1 comparison study. Two-axis classification:
%   axis-name carrier:    string-literal  vs  type-parameter  vs  value
%   check phase:          runtime          vs  compile-time

\TODO{landscape table; axes = ``where axis names live'' $\times$
``when they are checked''. Place einops, einx, Haliax, Penzai, xarray,
named PyTorch, tinygrad, and our system in this 2-D grid. Result: our
system uniquely occupies the (type-parameter, compile-time-when-known)
cell. Compress to ${<}1$ column of body text.}

\subsection{What remains open}
\label{sec:bg-gap}

The landscape leaves an explicit gap: \emph{compile-time axis checking
that does not abandon the runtime-shape case}. The static approaches
in the literature require either fully static shapes (no batch
flexibility) or full type erasure at the kernel boundary. Our design
target is to recover the latter without giving up the former.

% ===========================================================================
\section{Design 1: The \texorpdfstring{\udltex{}}{\_tex} Bridge}
\label{sec:udl}

% ~1.5 pages. The big idea: the formula IS the program. Show the
% parse, show what it expands to, show that it's a real C++ feature
% (consteval + UDL).

\paragraph{The interface.}
A contraction between an embedding matrix and a token-index vector
appears in code as:

\begin{lstlisting}[language=cppx]
auto y = "ij,j->i"_tex(E, x);
\end{lstlisting}

The literal \udltex{ij,j{->}i} is parsed \emph{at compile time} into a
structural representation of the contraction's axes, which becomes a
template parameter on the underlying kernel call.

\subsection{Parse pipeline}
\label{sec:udl-parse}

% The UDL operator hands the raw character pack to a consteval parser
% that produces a structural type: { lhs_axes, rhs_axes, output_axes }.
% Walk through the small grammar.

\TODO{walk-through of the consteval parser: grammar (one-line BNF),
the parsed AST type, how it becomes a non-type template parameter on
the kernel.}

\subsection{Why a UDL, not a string-typed API}
\label{sec:udl-why}

% Three reasons: (a) parse fires at compile time, errors land at the
% literal's source location, (b) the type system sees the parsed
% structure, not the string, (c) zero runtime overhead.

\TODO{three-paragraph contrast against \code{einops.rearrange("...")},
emphasising that the UDL pays the parse cost once at compile time and
the runtime side sees a typed structural value, not a string.}

\subsection{What the UDL is not}
\label{sec:udl-limits}

% Honest scope: it's not a full einops/einx replacement; broadcasts
% are explicit; the formula must be a literal (not constructible at
% runtime); ellipsis and partial application are deliberate non-goals.

The \udltex{} literal accepts a restricted Einstein notation: explicit
axes only, no ellipsis, no rank-polymorphic patterns. Patterns whose
shape varies at runtime are handled by the hybrid system in
\Cref{sec:hybrid} rather than by extending the literal grammar.

% ===========================================================================
\section{Design 2: Hybrid NTTP $\times$ DynamicShape}
\label{sec:hybrid}

% ~1.5 pages. Two compile-time paths and one runtime path; how they
% compose at the kernel boundary.

\subsection{The two paths}

A named tensor carries its axis information in one of two ways:

\begin{description}
  \item[Compile-time path] axis names \emph{and} sizes are encoded as
        non-type template parameters. The compiler refutes
        \udltex{ij,j{->}i} applied to an $(I,K)$ matrix; the error
        message names the offending axis.
  \item[Runtime path] axis names are encoded as template parameters
        (compile-time) but sizes live in a \code{DynamicShape} value
        (runtime). The compiler refutes axis-\emph{name} mismatches;
        size mismatches surface as a \code{ShapeError} at the kernel
        call site.
\end{description}

\subsection{Promotion and demotion}

% When a static-shaped tensor is fed into a code path that only knows
% the names (e.g., a generic kernel over a batch axis of unknown
% size), the static sizes "demote" transparently. The reverse
% promotion is explicit and checked.

\TODO{promotion/demotion rules with a one-page minimum example
(static-shaped activations flowing into a runtime-batched op).
Include the C++ snippet showing what the user writes vs what the
compiler instantiates.}

\subsection{Kernel boundary}

The contract of a kernel is a function template parameterised by the
axis structure (compile-time) accepting \code{DynamicShape} runtime
sizes when needed. \emph{The same kernel} services both paths; the
runtime-path overhead is one bounds check per call, not a separate
kernel implementation.

% ===========================================================================
\section{Design 3: Hexagonal-lite Kernel Backend}
\label{sec:hex-lite}

% ~1 page. The port; the three adapters; why this matters for MLSys
% specifically (substrate evolution, e.g., webgpu, is the active
% research area).

The named-axis layer should not know whether contractions are computed
by hand-rolled loops, an Eigen template instantiation, or a WGSL
kernel dispatched to a GPU. \emph{Hexagonal-lite} is a constrained
form of the Hexagonal architecture~\cite{adr0021}: a single port
(\code{KernelBackend}) with three concrete adapters.

\subsection{The port}

\TODO{one-page C++ listing of the \code{KernelBackend} concept (in the
C++20 \code{concept} sense) --- the minimal operations a substrate
must provide, with brief commentary on each requirement.}

\subsection{Substrates}

\begin{description}
  \item[Reference] portable naive nested-loop implementation; defines
        the semantics that the other substrates must match.
  \item[Eigen] BLAS-routed contractions for CPU performance baseline;
        loop fusion via Eigen's expression templates.
  \item[WebGPU via Dawn]~\cite{dawn,adr0014} GPU substrate compiled to
        WGSL; vcpkg-vendored Dawn to keep substrate provisioning
        reproducible in CI.
\end{description}

The choice of substrate is a CMake variable
(\code{-DTENSOR\_KERNEL\_BACKEND=\{reference,eigen,webgpu\}}) ---
the named-axis layer is unchanged across all three.

% ===========================================================================
\section{Evaluation}
\label{sec:eval}

% ~2 pages. Two evaluations:
%   §5.1 Static-catch benchmark (does the compile-time path actually
%        reject the axis-bug taxonomy?)
%   §5.2 Runtime case study (bundle adjustment; cross-backend perf;
%        comparison against einops + Haliax baselines).
%
% Both depend on B-stage bench harness. Numbers are [BENCH:...]
% placeholders for now.

\subsection{Static-catch benchmark}
\label{sec:eval-static}

% Methodology: an axis-bug benchmark suite (each entry is a small
% program known to be buggy). For each entry, score whether the
% library catches the bug at compile time, at first runtime call, or
% silently produces wrong output.

\BENCH{Static-catch benchmark suite: $\sim$50 known-buggy programs
spanning the taxonomy in \Cref{sec:bg-taxonomy}. Score each entry as
\{compile-time, runtime, silent\} for each of: \code{tensor} (ours),
\code{einops}, \code{einx}, \code{Haliax}, \code{Penzai}. Headline
table = library $\times$ \{caught at CT, caught at RT, silent\} with
percentages.}

\subsection{Bundle-adjustment case study}
\label{sec:eval-ba}

\BENCH{50-view $\times$ 200-landmark bundle-adjustment scene from a
public dataset (e.g., BAL). Three backends (reference / Eigen /
WebGPU). Metrics: time-to-residual, time-to-gradient, peak memory.
Comparison baselines: \code{einops.rearrange}-based Jacobian
assembly, \code{haliax.NamedArray} on JAX backend.}

\TODO{at minimum a single table with rows = backend, columns =
metric. Append a small figure showing scaling with view count if
space permits.}

\subsection{Threats to validity}

\TODO{honest discussion: benchmark size, single-author bias on the
axis-bug suite, hardware/dataset variability for the BA case study.}

% ===========================================================================
\section{Related Work}
\label{sec:related}

% ~1 page. Compressed M1 landscape report. Group by paradigm:
%   Named-string runtime: einops, einx
%   Named language-level (runtime): Haliax, Penzai, xarray, named torch
%   Tile/kernel-DSL adjacents: tinygrad, Halide, ATL
%   Theory: Chiang & Rush (TMLR)

\TODO{1-page synthesis from the M1 study \cite{landscape-comparison-2026}.
Each paragraph closes with what our system carries over and what it
diverges on. Notable: Chiang \& Rush \cite{chiang2024named} give the
notation; we give a type-system implementation of it.}

% ===========================================================================
\section{Discussion and Limitations}
\label{sec:discussion}

% ~0.7 pages. Where the design helps; where it does not; what comes
% next.

\paragraph{When to reach for the compile-time path.}
\TODO{small decision-aid: shapes known at site of use $\to$ CT path;
batch size variable but axes named at compile time $\to$ hybrid;
fully dynamic (e.g., loaded from disk) $\to$ runtime path with
upgrade-to-CT after the first call.}

\paragraph{Limitations.}
\TODO{honest enumeration:
(L1) UDL grammar deliberately small (no ellipsis, no rank polymorphism);
(L2) compile-time error messages are template-instantiation-deep ---
we provide diagnostic helpers but C++'s default error rendering is
unhelpful;
(L3) WebGPU substrate is in early adoption; portability beyond
Dawn-supported platforms is currently constrained.}

\paragraph{Future work.}
\TODO{2--3 sentences: extending the UDL grammar towards full einx
expressivity without losing the compile-time guarantee; auto-tuned
substrate selection; integration with Triton-class kernel languages.}

% ===========================================================================
\section{Conclusion}
\label{sec:conclusion}

Named-axis libraries should treat the host language's type system as
load-bearing infrastructure, not as a place to attach metadata. The
\code{tensor} library demonstrates that the static and dynamic worlds
compose: a compile-time UDL bridge, a hybrid NTTP $\times$
DynamicShape axis system, and a Hexagonal-lite kernel backend deliver
substrate-portable, axis-typed tensor programs in a mainstream
systems language. \TODO{one-sentence summary of the headline empirical
result once \BENCH{} numbers land.}

% ===========================================================================
\bibliographystyle{plain}
\bibliography{refs}

\end{document}
